\section{Introduction}
\label{sec:introduction}

Git histories turn software into a dynamical system: a repository moves
through states as commits apply edits.  The first \CodeWM{} preprint
studied whether a tiny JEPA-style model could learn this transition
process from code states and edit actions, and found a practical
stability barrier.  Under the usual fast EMA target update, training hit
a repeatable 700-step ceiling because the target encoder collapsed; a
near-static target encoder fixed that failure and enabled stable
15K-step latent edit prediction~\citep{akbulut2026collapse}.

This paper starts after that fix.  If a tiny JEPA can be stabilized, the
next question is not merely whether it can predict the next target
embedding.  The next question is where transition information lives
inside the model, why some seemingly natural prediction objectives fail
even when training is stable, and whether the final readout remains
discriminative enough for the prediction to mean anything.

The answer is geometric.  The previous six-loop encoder does not retain
transition signal uniformly across depth.  Early loop outputs contain a
measurable edit signal, but repeated application of the shared encoder
block progressively smooths that signal away.  By the final loop,
consecutive states are almost indistinguishable under the target
encoder, deltas are tiny relative to state vectors, and the representation
space is strongly anisotropic.  In that geometry, explicitly predicting
$z_{t+1}-z_t$ is the obvious experiment and the wrong substrate.

The positive result is deep supervision.  Instead of asking the final
representation to carry every signal, we train intermediate loop outputs
to remain predictive.  The current winning configuration uses a
three-loop encoder and predictive auxiliary losses on loops 1 and 2,
while retaining the near-static EMA target and looped predictor from the
stabilized \CodeWM{} recipe.  Deep supervision repairs the loop
intermediates, but it is not sufficient by itself: a high-lift
\code{loops3-auxpred} checkpoint still collapsed the final attention-pool
readout to roughly two effective dimensions.  Adding VICReg-style
variance and covariance regularization is the current Phase~9
breakthrough because it preserves the lift while restoring a
high-rank, discriminative readout.

\paragraph{Contributions.}
\begin{enumerate}
    \item \textbf{Direct state-space delta prediction is a negative
          result, by geometry.}  Stabilized six-loop \CodeWM{}
          compresses transition signal by roughly $12\times$ across
          loops (transition-to-state norm ratio $0.123 \to 0.010$) into
          an effective rank of $5/128$.  In that geometry, early-layer
          readouts, normalized delta losses, and residual predictors
          all fail to recover a competitive predictor.  The obvious
          objective is the wrong substrate; this is a reusable
          lesson for any tiny world-model builder.
    \item \textbf{Deep supervision for looped encoders.}  Predictive
          auxiliary losses on intermediate loop outputs preserve useful
          state representations and raise lift, but also expose that
          lift alone is not a sufficient success metric.
    \item \textbf{Readout-collapse diagnostics and VICReg repair.}
          Discriminability probes show that high lift can coexist with
          online effective rank $1.95/128$ and KNN@1 near random.  On a
          full-validation audit (N=$5000$), VICReg delivers a
          $\sim\!20\times$ improvement in readout discriminability
          ($14\times$ KNN@5, $25\times$ effective rank) over the
          pre-VICReg baseline, peaking at step 28K before monotonic
          overfit.
    \item \textbf{Head-to-head against code foundation
          models~\citep{guo2022unixcoder,feng2020codebert}.}  On a paper-grade retrieval
          eval (1000 queries over a 5000-commit gallery), the
          $1.3$M-parameter step-28K \CodeWM{} matches $125.9$M
          UnixCoder: a $+0.030$ MRR edge on moderate criteria and a
          statistical tie on the strictest action-cosine threshold, at
          $10\times$ lower CPU latency.  A smaller development-scale
          benchmark dominates CodeBERT ($124.6$M) at $21\times$ lower
          latency.
    \item \textbf{Evaluation corrections and bottleneck diagnosis.}
          A 128-sample gallery inflates absolute KNN numbers
          $\sim\!12\times$ relative to a 5000-sample gallery; we
          document the inflation explicitly.  Cross-seed runs reveal
          that the peak result is a single-seed optimum
          ($2\text{--}3\times$ KNN@5 variance), and that predictor
          effective rank ($5\text{--}6/128$) is the binding bottleneck
          independent of encoder health.
\end{enumerate}

\paragraph{Status.}
This is an active draft.  Phase~8 geometry and delta-prediction results
are treated as established within the current experiment log.  Phase~9
numbers are audited on the full validation set.  The foundation-model
head-to-head uses the peak step-28K checkpoint (seed~42, pre-fix) scaled
to $1000$ queries $\times$ $5000$ gallery.  A post-hoc seed bug was
identified: the step-28K peak benefited from a fortunate initialization.
Cross-seed runs (seeds 42, 43 post-fix) trail by $2\text{--}3\times$ on
KNN@5, and predictor collapse (effective rank $5\text{--}6/128$) is
present across \emph{all} seeds.  Cross-seed variance and the predictor
bottleneck are reported honestly in \S\ref{subsec:cross_seed}.  A
CodeT5+ retry with the correct encoder API and \KernelWM{} transfer
remain in progress.
